

\subsubsection*{Replace the assessment while retaining the payments}
After further company activity, the new period's cumulative tax calculation may
change. Its recorded payments do not change with it. A replacement must therefore
start from the actual continued book and distinguish what was paid against the
carried liability from what was paid against the current period's assessment.
Using current expense less all tax payments would wrongly treat a payment of
last year's tax as a reduction in this year's expense.

Let $C_0$ be recorded current-period tax expense, $C_1$ the replacement cumulative
amount, $S$ the actual payments allocated to that current vintage, and $O$ the
remaining carried liability. The replacement records only the expense difference:
\[
  \Delta C=C_1-C_0,\qquad
  P_{\mathrm{current}}=C_1-S,\qquad
  P_{\mathrm{total}}=O+C_1-S.
\]
The journal debits current-tax expense by $\Delta C$ and credits payable by the
same signed amount. Cash and the carried liability remain unchanged. The current
vintage retains the identity referenced by earlier payments; its old assessment
and settlement record remain available alongside the replacement's provenance.

Continue the split-payment example above. Recorded current expense is 5,
current-vintage payments are 2, and carried payable is $15/4$. If the replacement
cumulative tax is 7, the expense increment is 2, current payable becomes 5 and
total payable becomes $35/4$. Cash stays at $715/2$ and recorded profit falls
from $-17$ to $-19$. Assets remain $1367/2$, while liabilities become $475/4$
and equity including profit becomes $2259/4$. If cumulative tax instead falls
to 3, the expense reduction is 2 and total payable becomes $19/4$.

A positive aggregate liability can conceal a different representation problem.
If cumulative current tax falls to 1 after payments of 2, the current vintage
would have a credit of 1 even though the carried liability leaves a positive
aggregate of $11/4$. The present positive-payable model refuses that replacement
and preserves the preceding book. A credit or refund needs a specified account
and treatment; it cannot be silently netted against the older obligation.

Losses are also recalculated cumulatively. Suppose the carried opening contains
an available loss of 19 and the first assessment uses all 19 against taxable
profit of 20. At a declared rate of one quarter, current tax is $1/4$. A
replacement using the same cumulative profit and deduction again yields $1/4$:
it restarts from the original opening inventory, rather than using the provisional
closing balance of zero or recording a second consumption of 19. Previous
provisional calculations remain in the history. If that loss expires before the
new assessment's usage date, the replacement must apply the stated expiry rule;
it cannot retain the earlier deduction automatically.

The cumulative interval must cover every completed event and cannot shorten the
previous assessment. A correction at the same endpoint is possible when no later
event lies outside it. Native return adjustments enter the replacement exactly
once; a worked return sequence has pretax income 6 and an adjustment of $-24$,
producing a new loss of 18 without an automatic deferred asset. Fresh deferred
recognition, another continuation or repeated replacement requires a compatible
method. Source admission, complete fiscal coverage and annual acceptance remain
unresolved.
